\color{unimain}{\subsection{Сигнализация}}
\color{black}

\begin{enumerate}

\item
Устройство формирует выходное воздействие на сигнализацию \textit{«Сигнализация: Срабатывание фикс»} и \textit{«Сигнализация: Срабатывание»} по схеме «ИЛИ» в следующих случаях:
\begin{itemize}
\item при срабатывании дифференциальной токовой защиты;
\item при срабатывании резервных токовых защит КО, ОУ и ошиновки НН;
\item при срабатывании КИ НН;
\item при срабатывании внешних защит.
\end{itemize}

\item
Устройство формирует выходное воздействие на сигнализацию \textit{«Сигнализация: Неисправность фикс»} и \textit{«Сигнализация: Неисправность»} по схеме «ИЛИ» в следующих случаях:
\begin{itemize}
\item при неисправности цепей тока дифференциальной токовой защиты;
\item при неисправности цепей оперативного тока;
\item Потеря GOOSE.
\end{itemize}

\item
Алгоритм работы логики сигнализации приведен на рисунке \ref{pict_ls:LS}. 

\medskip
\begin{figure}[H]
\centering
\includegraphics[width=1\textwidth,height=1\textheight,keepaspectratio]{\fbpath/11.02 УШР/Сигнализация/img/LS.pdf}
\captionof{figure}{Функциональная схема логики сигнализации}\label{pict_ls:LS}
\end{figure}

\end{enumerate}